\section*{摘要}
药物母体采样不均与条件测量重复，使溶解度模型需要同时处理结构差异和温度响应的跨药物迁移。本文提出SOLAR（溶曜）母体均衡结构--温度分解框架：SOLAR-H以树集成学习结构基线，通过组合内回归估计温度响应，并由分组验证选择响应增益；SOLAR-N以双分子编码器和温度差分监督实现联合学习。基于1904条记录，在1507条开发记录、128个母体上开展嵌套分组验证，比较12种方法。SOLAR-H的母体等权MAE为0.7164，相对ExtraTrees的配对误差差为-0.0305，95\%区间为[-0.0693, 0.0086]。固定配置消融中，将温度增益设为1使MAE增加0.0987，支持通过收缩控制条件响应的迁移误差。特征归因与组合固定效应分析表明，结构信息主导当前预测，而组合内温度关联与跨药物预测收益具有不同含义。该框架为小样本条件溶解度建模提供了兼顾母体均衡与响应调节的实现方式。

\noindent\textbf{关键词：}药物溶解度；母体均衡；结构交互；温度收缩；分组验证
